\section{Proposed Approach}
\label{sec:proposal}
% on est sur de cette section ? 
% perso je l'aurais enlevé pour avoir plus de place par expériences.
% A la place on peut dire un mot sur les modèles utilisés LLama 3.1 8B en étude principale avec Qwen 3.8 27B en extension
% le dataset de Hahami
% la manière de compute les steering vectors comme dans le hahami et comment on a appliqué le dropout bruit gaussien et vecteur aléao
Our framework addresses three questions: what determines perturbation
detectability, which alternative explanations survive behavioral controls,
and how precisely a model can report information about an internal
intervention. The following experiments examine these questions separately;
the behavioral hypotheses are not inferred from the calibration procedure
used to define intervention strength.

\subsection{Perturbation Type, Dose Matching, and Depth}
\label{sec:localization-approach}

We study Llama-3.1-8B-Instruct and intervene at the outputs of its 32 decoder
blocks, with Qwen3.8-27B and its 64 blocks as an extension. The two models
differ in how much a forced choice can reveal: Llama holds a large standing
preference for one answer label, so its decision rarely moves, whereas Qwen's
preference is small enough that its decision does move. Comparing them
separates what a perturbation does from what a discrete report exposes. The main design compares four perturbation families: concept-derived
directions, fixed isotropic random directions, random directions renewed at
each targeted token, and dropout. Concept-to-random comparisons test whether
detectability depends on the structured content of the injected direction.
Comparing fixed and renewed random directions isolates the effect of
cross-token directional coherence, while dropout provides a qualitatively
different, activation-dependent control.

For additive interventions, we apply
\begin{equation}
h'_{\ell,t}=h_{\ell,t}+\alpha v,
\qquad
\|v\|_2=1,
\label{eq:intervention}
\end{equation}
where $\alpha$ is the Euclidean displacement applied at each targeted token.
Equal values of $\alpha$, however, need not represent equally unusual
interventions because natural activations vary differently across layers
and directions. We therefore introduce a layer- and direction-specific
scale $s(\ell,v)$ to complement raw-amplitude matching.

For each unit direction $v$, we project unperturbed activations collected at
the future intervention site:
\begin{equation}
p_i(\ell,v)=\langle h_{\ell,i},v\rangle.
\label{eq:projection}
\end{equation}
Our primary scale is the sample standard deviation of these projections,
\begin{equation}
s_{\mathrm{SD}}(\ell,v)
=
\sqrt{
\frac{1}{N-1}
\sum_{i=1}^{N}
\left(
p_i(\ell,v)-\bar{p}(\ell,v)
\right)^2
}.
\label{eq:scale}
\end{equation}
It converts the raw amplitude into a standardized dose:
\begin{equation}
z=\frac{\alpha}{s(\ell,v)},
\qquad
\alpha=z\,s(\ell,v).
\label{eq:dose}
\end{equation}
Matching $\alpha$ preserves the absolute displacement, whereas matching $z$
controls the displacement relative to the natural variability of the
corresponding layer and direction. These two parameterizations therefore
test whether apparent differences between perturbation families persist
after accounting for activation geometry.

Scales are estimated separately for each $(\ell,v)$ pair from unperturbed
activations recorded at the same block output and eligible token positions
as the subsequent intervention. The sample SD is the primary definition of
$s(\ell,v)$, while the corrected median absolute deviation provides a robust
alternative parameterization. Uncertainty is estimated by resampling
sentences while retaining all eligible tokens belonging to each sampled
sentence. Because activation statistics can depend on prompt structure,
the scales used for behavioral comparisons are computed in the
corresponding presentation context.

Two-alternative forced-choice localization serves as the main behavioral
measurement instrument. Each prompt contains two sentences, one of which is
perturbed, and the model must identify the targeted sentence. Sentence
order, answer labels, and intervention target are counterbalanced so that
the measured effect follows the perturbed sentence rather than a particular
label or physical position.

Accuracy alone cannot carry this measurement, because a perturbation that
carries no information about its target still displaces the answer logits
away from the model's default label, and that displacement scores as a
success on every trial whose target happens to carry the other label. We
therefore report a paired contrast between two trials that share every factor
except which sentence was perturbed,
\begin{equation}
S = \tfrac{1}{2}\left[L(\text{A targeted})-L(\text{B targeted})\right],
\label{eq:paired}
\end{equation}
where $L$ is the sham-subtracted difference between the two answer logits.
Any displacement independent of which sentence was targeted cancels exactly,
and a permutation of the target assignment within each pair provides the null.

Because a concept direction differs from a random one in both its orientation
and its coordinate statistics, we add a \emph{scrambled concept} control: a
coordinate permutation of a concept vector, which preserves its norm and the
multiset of its coordinates while destroying alignment with any feature
direction. Comparing a concept direction against its own permutation at
matched amplitude isolates orientation from magnitude.

Dose sweeps are evaluated under both raw-amplitude and standardized
matching. The primary comparisons ask whether psychometric responses differ
between perturbation families at matched $\alpha$, and whether these
differences remain at matched $z$. Variation across layers, concepts, and
random directions is retained rather than reduced to a single family
average. A 75\% threshold is reported only when it is bracketed by the tested
dose range and supported by the observed curve shape; nonmonotonic behavior
at high doses is treated as a result rather than extrapolated away. The
executed sweeps, dose grids, trial counts and sham design are listed in
\cref{app:loc-design}.

\subsection{A Semantic Control Task}
\label{sec:degradation-approach}

Localization measures what a perturbation leaves legible, not what it costs.
A second forced choice therefore asks nothing about the intervention: the
model reads one perturbed sentence and answers, with one of two letters,
whether it expresses a queried concept. That concept is never the injected
one, so a targeted injection cannot supply the answer, and both letter
assignments are run on the same sentence, so a letter preference stays
separable from a classification. Each trial is scored by the sign and the size
of the logit difference between the two letters against the sham response of
the same item, and the blocks, families and doses are those of the
localization sweep, so every cell pairs with a localization cell.

\subsection{Multi-injection for model perturbation}
\label{sec:multi-injection-approach}

A complementary question is whether reports remain informative when several concept-derived interventions are present at once: can the model detect, identify and locate those interventions, rather than only report something changed at all?

We extend the additive intervention of \cref{eq:dose} to $n$ simultaneous injections, $v_i, i\in\{0,\dots,n\}$ , each concept direction $v_i$ applied at its own layer $\ell_i$ and sustained through generation.
The model's behaviour under multi-intervention is evaluated through three tasks.
\emph{Counting} injects $k\in\{0,\dots,4\}$ concept directions at distinct layers and requires a single forced digit response, with $k{=}0$ serving as the true sham.
\emph{Identification} injects two concepts and asks the model to name what was injected, scored both by forced choice among labeled candidates and by free response.
\emph{Relative-depth ordering} injects two concepts at two layers and asks which was injected at the shallower layer, presented as a randomized \texttt{A}/\texttt{B} choice so that label assignment is decorrelated from true depth. Details of the different task prompts can be found in Appendix.
All tasks draw from a pool of ten concept directions (five simple, five complex) applied across all 32 decoder blocks, with a coherence filter excluding degenerate or non-answer generations before scoring.
For this program, we use raw amplitude $\alpha$ instead of the standardized dose $z$ described in \cref{eq:dose}; this is a scope choice distinct from the raw-versus-standardized comparison detailed before.

A discrete decision can hide a default-answer preference unrelated to the injection, so each task pairs its accuracy metric with a control that isolates the injection's own contribution.
A \emph{bias-adjusted logit contrast} compares the first-token logit for the true answer against the model's default (e.g.\ $\mathrm{logit}(k)-\mathrm{logit}(\text{``1''})$ for counting), read directly from the generation call so it is available even for incoherent trials.
For ordering, a \emph{paired-sham} control repeats each injected trial with zero injection on the same prompt, so that $\text{contrast}_{\text{injected}}-\text{contrast}_{\text{sham}}$ isolates the injection's effect from the prompt's own baseline preference.
